\chapter{RELATED WORK}

This chapter surveys the relevant literature across four interconnected areas that inform CA-DQStream's design. Section 2.1 reviews streaming data quality monitoring frameworks. Section 2.2 covers concept drift handling in ML systems. Section 2.3 addresses context-aware anomaly detection. Section 2.4 examines memory-based streaming anomaly detection. Section 2.5 summarizes the chapter.

% ================================================
\section{Streaming Data Quality Monitoring}
% ================================================

Recent work has extended data quality monitoring to streaming contexts, addressing some of the limitations of batch-oriented approaches.

\textbf{Stream DaQ}, developed by Papastergios and Gounaris, employs continuous windowing mechanisms with dynamic constraint adaptation for real-time data quality monitoring \citep{papastergios2025stream}. The framework validates streaming data against user-defined constraints that adapt as data distributions evolve, demonstrating feasibility of online constraint evaluation. However, Stream DaQ relies exclusively on rule-based validation without ML-based anomaly detection.

\textbf{Grab's Coban Platform} uses FlinkSQL for real-time semantic validation through Kafka stream contracts \citep{grab_coban2025}. Deployed in early 2025, Coban monitors over 100 critical Kafka topics and enables immediate identification and halting of invalid data propagation. The system employs LLM-based semantic test rule recommendation to suggest domain-specific constraints from provided schemas and anonymized sample data, reducing the cognitive burden of manual rule authoring. Coban demonstrates that FlinkSQL-based test execution is production-viable for enterprise-scale streaming DQ workloads.

\textbf{Confluent's streaming quality framework} provides a six-step validation pipeline architecture that has shaped industry practice: (1) ingest data into Kafka topics, (2) enforce schema validation via Schema Registry at ingestion, (3) apply business rules via CEL-based data contracts, (4) route invalid records to quarantine for review, (5) monitor quality metrics continuously, and (6) trigger alerts when thresholds are violated \citep{confluent_2025}. The shift-left paradigm moves validation to ingestion time, stopping bad data before it spreads. Siemens Healthineers processes 8 million messages daily via Confluent, catching manufacturing defects instantly \citep{siemens_confluent2024}. These deployments validate that streaming DQ monitoring is operationally essential, but they focus on syntactic validation and semantic rule checking rather than ML-based anomaly detection for multivariate quality violations.

\textbf{Rapidash} addresses complex multi-record violations using spatial indexing for streaming contexts, extending the scope of streaming DQ to multi-record consistency checks.

The streaming DQ landscape reveals significant gaps. No existing framework combines ML-based multivariate anomaly detection with context-aware thresholding, sequential pipeline architecture, and multi-strategy drift adaptation in a unified system. Production deployments of streaming platforms provide additional architectural insights and operational benchmarks.

\textbf{Grab's Signals Marketplace} implements a data mesh architecture for a ride-hailing, food delivery, and financial services platform serving 800 cities across 8 Southeast Asian countries \citep{grab_streaming_dq}. The platform introduces data certification with SLA guarantees, where certified data assets carry formal schema contracts and completeness/freshness commitments. Deprecated tables increased 400\% year-over-year after data mesh adoption, and 75\% of Grab queries now hit certified assets, demonstrating the operational impact of formal DQ governance.

\textbf{Grab's Trailblazer CDC Pipeline} demonstrates Change Data Capture via MySQL binlogs using Debezium on Kafka Connect, with Spark Structured Streaming persisting to a petabyte-scale data lake \citep{grab_streaming_dq}. The architecture (binlog $\rightarrow$ Debezium $\rightarrow$ Kafka $\rightarrow$ Spark $\rightarrow$ data lake) serves as a reference topology for streaming pipeline instrumentation. A DASH (Data Auditor as a Service) component performs hourly checks comparing ID counts between source databases and the streaming layer.

\textbf{Grab's FlinkSQL Interactive Platform} demonstrates a three-layer architecture (Compute, Integration, Query) for deploying streaming pipelines at scale \citep{grab_streaming_dq}. Cold start reduced from 5 minutes to 1 minute through shared cluster strategies. Full pipeline deployment completes in under 10 minutes. The three-layer separation of concerns (infrastructure vs. integration logic vs. query interface) informs the modular design of CA-DQStream's processing pipeline.

\textbf{Grab's AutoMQ Migration} demonstrates shared-storage Kafka architecture achieving 3$\times$ throughput increase and 3$\times$ cost efficiency improvement over traditional Kafka deployments \citep{grab_streaming_dq}. Partition reassignment time dropped from 6 hours to under 1 minute through metadata-only synchronization.

\textbf{Vimeo} reports that batch validation introduced a one-day delay before insights reached analytics teams, limiting real-time decision-making and quick pivots after launches and campaigns \citep{confluent_2025}. This operational cost of batch-oriented DQ validation motivates the shift to streaming-first quality monitoring. The Confluent 2025 Data Streaming Report surveys global streaming adoption, finding that only approximately 1\% of companies have reached the highest streaming maturity level where data streams are treated as managed products with strategic quality guarantees \citep{confluent_2025_streaming_report}.

\autoref{tab:dq_framework_comparison_related} provides a qualitative comparison of CA-DQStream against existing frameworks.

\begin{longtable}{|m{2.8cm}|m{2cm}|m{1.6cm}|m{1.6cm}|m{1.6cm}|m{1.8cm}|}
\hline
\textbf{Framework} & \textbf{Type} & \textbf{ML Detection} & \textbf{Streaming} & \textbf{Context-Aware} & \textbf{Drift Adapt} \\ \hline
Great Expectations & Batch & \ding{55} & \ding{55} & \ding{55} & \ding{55} \\ \hline
Soda Core & Batch & \ding{55} & \ding{55} & \ding{55} & \ding{55} \\ \hline
AWS Deequ & Batch & \ding{55} & \ding{55} & \ding{55} & \ding{55} \\ \hline
Stream DaQ & Streaming & \ding{55} & \checkmark & \ding{55} & \ding{55} \\ \hline
IForestASD & Streaming & \checkmark & \checkmark & \ding{55} & \checkmark \\ \hline
CA-DQStream & Streaming & \checkmark & \checkmark & \checkmark & \checkmark \\ \hline
\caption{Qualitative comparison of CA-DQStream against existing DQ frameworks.}
\label{tab:dq_framework_comparison_related}
\end{longtable}

% ================================================
\section{Concept Drift Handling in ML Systems}
% ================================================

Concept drift handling is a well-studied problem in machine learning, with a rich literature on detection methods, adaptation strategies, and theoretical analyses of learning under non-stationary distributions \citep{gama2014survey, lu2018learning}. The detect-then-retrain paradigm underpins most traditional approaches: when drift is detected via statistical monitoring, the system triggers full model retraining. This is conceptually straightforward but computationally expensive and introduces significant adaptation latency during which the stale model continues to produce unreliable results.

\textbf{Baier et al.} proposed a switching scheme combining periodic retraining with triggered incremental updates, demonstrating that NYC taxi demand is particularly susceptible to incremental drift due to weather, events, and economic factors \citep{baier2020switching}. Their scheme compared ADWIN, STEPD, and HDDDM, finding that any drift detection strategy outperforms no strategy, though differences between specific strategies are significant only under certain drift characteristics.

\textbf{DTD (Dynamic Threshold Determination)} formally proves that fixed thresholds are suboptimal and demonstrates that DTD-enhanced HDDM-W achieves 58.31\% accuracy versus 48.64\% with fixed thresholds on the Airline dataset \citep{dtd2026}. The three theorems establish: (1) perfect detection (zero false alarms, zero delay) may not be optimal; (2) no single fixed threshold is universally optimal; and (3) a dynamic threshold strategy strictly outperforms any fixed threshold.

\textbf{Multi-strategy drift adaptation} represents an emerging frontier. Rather than applying a uniform response to all drift events, these approaches classify drift type and severity, then select appropriate strategies. METER (MEta-learning-based adapTER) uses an MLP hypernetwork trained on historical drift patterns to predict optimal adaptation parameters from current meta-metrics \citep{meter_2023}, providing the algorithmic foundation for moderate drift adaptation where model parameters shift but overall structure remains. The SCAR benchmark provides the most comprehensive evaluation of streaming anomaly detection under concept drift, testing 9 streaming algorithms and 4 adapted static baselines across 76 synthesized and 74 real-world datasets \citep{scar2025benchmark}. Key findings include: (1) no single algorithm statistically dominates; (2) adapted static baselines (iForest, LOF, iNNE, IDK) with reconstruction strategies rank among top performers; (3) model update strategy matters more than algorithm choice; and (4) memory-augmented methods may underperform on high-dimensional data. CA-DQStream's Intelligent Evolution Controller extends multi-strategy adaptation with four complementary strategies (Continuous Evolution, Switching, METER, Spatial Tracking) selected dynamically based on drift characteristics.

% ================================================
\section{Context-Aware Anomaly Detection}
% ================================================

Context-aware anomaly detection conditions the notion of anomalousness on the context in which observations occur. Foundational work by Dey established the theoretical framework for context-awareness in computing systems \citep{dey2001understanding}, while subsequent work has explored context-aware approaches in sensor networks, financial fraud detection, and healthcare monitoring. In data quality, the key insight is that threshold values for detecting anomalies must be calibrated to the specific context of each observation, not applied uniformly across heterogeneous data streams. Systematic literature reviews have identified context as a critical but under-defined concept in DQ management \citep{contextdq2022}. DTD formally proves that no single fixed threshold can be universally optimal across all data segments \citep{dtd2026}.

\textbf{Multi-dimensional thresholding approaches} partition the feature space into regions and compute separate detection parameters for each region. These can be implemented through decision trees (leaf nodes as context-specific thresholds), mixture models (components as sub-populations), or explicit context matrices. \textbf{ReAD (Regional Anomaly Detection)} uses dynamic partitioning based on Delaunay triangulation, computing thresholds that adapt to local data density through a relative divergence metric \citep{read2020}. The key insight is that the same absolute value is interpreted differently depending on neighboring regions: high ridership in a city center is normal, but the same value in a suburban area is anomalous. This relative interpretation mirrors the context-aware principle that quality thresholds must be calibrated to local conditions.

\textbf{Trajectory-based approaches} use expected behavior as a reference for anomaly detection. \textbf{TAPS (Taxi Anomaly detection via Prediction-based detection System)} employs an offline-training / online-detection architecture where trajectory prediction models trained on normal behavior are deployed for real-time detection \citep{taps2021}; deviations exceeding learned thresholds are flagged as anomalous. \textbf{MTRI (Multi-scale Temporal and Road Network Interaction Anomaly Detection)} captures multi-scale temporal features combining contrastive learning-based conditional diffusion models with heterogeneous graph modeling of road network interactions, achieving AUC-ROC > 0.85 on real-world GPS trajectory data \citep{mtri2026}. The method demonstrates that multi-scale temporal modeling (short-term speed patterns vs. long-term route deviations) is essential for comprehensive anomaly detection.

\textbf{Behavior-profile-based approaches} maintain profile representations of normal behavior. \textbf{BeSTAD (Behavior-Aware Spatio-Temporal Anomaly Detection)} detects anomalies in human mobility data by learning context-specific behavioral norms across six dimensions: Jensen-Shannon divergence of cluster distributions, dominant mode change, new pattern emergence, transition structure change, entropy change, and frequency change \citep{bestad2025}. BeSTAD achieves AUROC = 0.775 compared to 0.586 for prior context-aware methods.

\textbf{Interpretable approaches} such as \textbf{ICAD (Interpretable Component-wise Anomaly Detection)} enable anomaly scoring at the individual feature level, decomposing anomaly scores into spatial versus temporal components \citep{icad2025}. This supports explainability by attributing anomalies to specific behavioral dimensions rather than holistic scores, critical for DQ systems where operators need to understand why a record was flagged.

The primary limitation across all existing approaches is reliance on manual threshold specification: practitioners must manually define thresholds for each context cell, leading to scalability challenges as context dimensions grow. No existing framework provides automated threshold computation from training data combined with intelligent fallback for sparse contexts and drift-aware recalibration.

\subsection{Anomaly Detection Methods and Benchmarks}

A comprehensive benchmark literature provides the methodological foundation for evaluating streaming anomaly detection systems.

\textbf{ADBench} evaluates 30 anomaly detection algorithms across 57 datasets, demonstrating that no unsupervised algorithm statistically dominates across all datasets \citep{adbench2022}. Key findings: (1) semi-supervised methods achieve median AUC-ROC of 80.95\% with only 5\% labeled anomalies versus 60.84\% for fully supervised; (2) ensemble tree methods (XGBoost, CatBoost) are competitive; (3) deep learning unsupervised methods (DeepSVDD, DAGMM) underperform shallow methods without label guidance; and (4) HBOS, COPOD, ECOD, and NB are the fastest methods, providing guidance for real-time scoring layers.

\textbf{Statistical methods} provide the speed-accuracy trade-off needed for real-time operation. \textbf{HBOS (Histogram-based Outlier Score)} constructs histograms for each feature and scores instances based on their deviation from bin frequencies, treating features independently \citep{adbench2022}. \textbf{COPOD (Copula-based Outlier Detection)} models feature dependencies through empirical copulas, providing stronger performance on correlated features while remaining fast \citep{adbench2022}. \textbf{ECOD (Empirical CDF-based Outlier Detection)} uses empirical cumulative distribution functions for faster computation and strong performance on mixed-type data \citep{adbench2022}. Together, these methods provide a fast first-pass detector layer that CA-DQStream's context-aware scoring can refine.

\textbf{Isolation Forest variants} address specific limitations of the original algorithm. \textbf{K-Means-based IF (K-IF)} replaces binary tree splits with multi-branch splits determined by K-Means clustering, demonstrating significantly improved anomaly ranking differentiation on the NYC Taxi dataset (3.38M records): 12,500 distinct rank values versus approximately 1,750 for standard IF, with 7$\times$ faster execution on geographic-only data \citep{karczmarek2020kmeans}. Critically, K-IF assigns zero scores to records with missing/invalid values in 99.98\% of cases versus 99.56\% for standard IF, directly relevant to the completeness validation challenge in streaming taxi data. \textbf{RIFIFI (Revised IF to Identify Fraud and data Inner structure)} modifies split selection based on density, generating Isolation Nodes, Dense Nodes, and Depth-Limit Nodes \citep{yepmo2024rififi}. RIFIFI achieves mean AUC of 0.88 versus 0.85 for standard IF, with particularly strong improvements on local anomalies in varying-density regions (+0.233 AUC on Mammography, +0.097 on Wood). This demonstrates that density-aware splitting improves local anomaly detection, relevant for distinguishing anomalies in taxi-dense versus taxi-sparse zones.

\textbf{IForestASD (Isolation Forest for Streaming Data)} adapts IF for streaming contexts through a sliding-window approach: maintaining a window of recent observations, rebuilding the forest when the anomaly rate exceeds a threshold, and scoring new instances against existing trees \citep{togbe2020iforestasd}. On the Shuttle dataset (7.15\% anomaly rate), IForestASD achieves F1 of 0.64--0.80 versus 0.13--0.17 for Half-Space Trees (4--5$\times$ improvement). On SMTP (0.03\% anomaly rate), IForestASD detects anomalies where Half-Space Trees completely fail (F1 of 0.34--0.40 versus 0). The method uses approximately 20$\times$ less memory than Half-Space Trees. Recommended window size is approximately 500 for best F1, with a trade-off between scoring speed and detection accuracy. IForestASD provides the most direct template for streaming IF adaptation in CA-DQStream's architecture.

\textbf{The Numenta Anomaly Benchmark (NAB)} provides labeled NYC Taxi data with 10,320 data points at 30-minute intervals containing five ground-truth anomalies: NYC Marathon, Thanksgiving, Christmas Day, New Year's Day, and the January 2015 Blizzard \citep{nab2020}. NAB establishes the benchmark protocol for evaluating anomaly detection on urban mobility data.

\textbf{Self-supervised anomaly detection} has emerged as state-of-the-art by learning representations through proxy tasks without requiring labeled data \citep{hojjati2025selfsupervised}. The survey evaluates 180+ recent studies across image, time-series, tabular, and graph data, demonstrating that self-supervised methods significantly outperform classical approaches (KDE, OC-SVM, Isolation Forest) on complex data. Relevant proxy tasks include predicting time-of-day or day-of-week from partial features, which could enrich CA-DQStream's contextual representations.

\textbf{Time-series anomaly detection} surveys across 20 methods demonstrate that Isolation Forest consistently performs well across datasets with minimal hyperparameter sensitivity \citep{braei2020timeseries}. Key findings include: (1) no universal winner across all data characteristics; (2) window size is critical (too small misses patterns, too large introduces noise); and (3) threshold selection (typically 3$\sigma$ or percentile-based) significantly impacts results.

% ================================================
\section{Memory-Based Streaming Anomaly Detection}
% ================================================

Memory-based streaming anomaly detection methods maintain explicit storage of normal patterns that can be queried at inference time to score new observations.

\textbf{MStream} extends memory-based detection with dynamic graph-based clustering, maintaining cluster structure alongside individual pattern representations \citep{siddiqui2019mstream}. The memory module stores real-valued vectors of encoded normal patterns, queried via $K$-nearest neighbors to compute anomaly scores. Memory poisoning prevention is achieved through a $K$-nearest neighbor discounting mechanism and an update threshold that only permits instances with low anomaly scores to modify the memory. The River library implements memory-augmented streaming anomaly detection variants, with configurable memory replacement strategies (FIFO, LRU, Random) and a grace period mechanism for encoder bootstrapping \citep{montiel2021river}.

The \textbf{SCAR benchmark} evaluates additional streaming algorithms: \textbf{iForestASD} (monitoring outlier rate changes in sliding windows), \textbf{ARCUS} (comparing anomaly score distributions between batches), \textbf{RRCF} (Robust Random Cut Forest with collaborative dispersion scoring), \textbf{LODA} (Lightweight Online Detector of Anomalies), \textbf{xStream} (sequential hypothesis testing on streaming data), and \textbf{STORM} (Self-Tuning Online Random Forest with fading factors) \citep{scar2025benchmark}. Key findings: (1) adapted static baselines (iForest, LOF) ranked in the top four among 13 streaming algorithms; (2) model update strategy matters more than algorithm choice; and (3) no single algorithm statistically dominates.

The gap in memory-based streaming anomaly detection is the lack of integration with context-aware mechanisms. Existing memory-augmented methods maintain a single unified memory for all observations, without conditioning memory content or scoring on context. This uniform treatment leads to context collapse, where memory captures average behavior rather than context-specific normal patterns. More fundamentally, no existing framework combines memory-based adaptation with multi-strategy drift monitoring.

CA-DQStream addresses these limitations through four innovations: (1) 4D context-aware thresholding that conditions anomaly detection on multi-dimensional context; (2) sequential multi-stage validation pipeline; (3) multi-strategy drift adaptation via the Intelligent Evolution Controller; and (4) validated performance on 72 million records spanning 26 months, demonstrating production-grade reliability.

% ================================================
\section{Chapter Summary}
% ================================================

This chapter situates CA-DQStream within the broader research landscape across seven interconnected areas. Batch-oriented DQ frameworks (Great Expectations, Soda Core, AWS Deequ) provide the operational maturity of data validation but cannot operate in streaming contexts. Streaming DQ monitoring frameworks (Stream DaQ, Coban, Confluent) validate the architectural patterns (schema enforcement, CEL-based contracts, DLQ routing) but lack ML-based anomaly detection and drift adaptation. Industrial case studies (Grab Signals Marketplace, Trailblazer, FlinkSQL Platform, AutoMQ) provide architectural templates and operational benchmarks for production deployment. Concept drift literature provides detection methods (ADWIN, ADWIN-U) and adaptation strategies (Baier switching scheme, SCAR multi-strategy evaluation) but has not been integrated into DQ monitoring contexts. Context-aware anomaly detection (ReAD, BeSTAD, ICAD, DTD) demonstrates that threshold calibration must be context-dependent but relies on manual specification. Anomaly detection benchmarks (ADBench, SCAR, NAB, IForestASD, K-Means IF, RIFIFI) establish rigorous evaluation methodology and identify HBOS/COPOD/ECOD as fast first-pass detectors. Memory-based streaming methods provide principled mechanisms for incorporating new patterns without forgetting but lack context-awareness.

The research gaps identified across these areas guide CA-DQStream's four core contributions: no existing framework combines context-aware thresholding with memory-augmented anomaly detection, sequential pipeline architecture, and multi-strategy drift adaptation in a unified streaming DQ system. Chapter 4 presents the CA-DQStream architecture, demonstrating how these four capabilities are integrated.
